
\section{Polynomial-time reductions}\label{flow: sec: reductions}\doHeader{Reductions}
So far our primary strategy for solving a given computational problem
has been to design an algorithm from scratch.
To this end, we've studied basic algorithm-design techniques
such as divide-and-conquer, dynamic programming, and greedy algorithms.
In the remainder of these notes, we consider a different approach:
solving the problem that you want to solve by \emph{reducing} it to a different problem,
one that already has a known algorithm.
More broadly, our interest is in how different computational problems relate to each other.
We start with some necessary formal definitions,
then illustrate the concept by describing how to reduce some new problems to \prob{Max Flow}.

\paragraph{Decision problems (definition).} 
When studying reductions, for simplicity and for historical reasons,
we usually cast the problems of interest as \emph{decision problems}---problems with yes/no answers.
That is, for any input, the output is either `\textsf{yes}' or `\textsf{no}'.
Most computational problems can easily be recast in this way.
For example, here is the decision version of \prob{Max Flow}:

\begin{mylist}
\item[] \prob{Max Flow} (decision version) \dotfill \em
\item[\bf input:] a flow network (digraph $G=(V,E)$ with edge capacities, source $s$ and sink $t$), and a $k\ge 0$
\item[\bf output:] `\textsf{yes}' if $G$ has an $s$-$t$ flow with value at least $k$, and otherwise `\textsf{no}'
\end{mylist}

\paragraph{Polynomial-time (definition).}
The first question we consider when approaching a new computational problem
is whether the problem has an algorithm that runs in \emph{polynomial time}---that is, for every possible input, the algorithm computes a correct output
in time bounded by some polynomial in the size of the input.
The degree of the polynomial can be large, as long as it's constant.
(For example, $O(n^{1000})$ is polynomial time.  $O(n\log n)$ is also polynomial time,
as $O(n\log n) = O(n^2)$.)

\paragraph{Polynomial-time reductions (definition).}
Now we can formally define a \emph{polynomial-time reduction}.
Let $A$ and $B$ be any two decision problems.
A \emph{reduction} from $A$ to $B$ is a function $r$
that maps each input $x$ for $A$
to a corresponding input, $r(x)$, for $B$.
The reduction is \emph{correct} if it has the following property:
\begin{quote}\em
  For every possible input $x$ to $A$,
  \\\hspace*{1em}the answer to $x$ is \textsf{yes} if and only
if the answer to $r(x)$ (as an instance of $B$) is \textsf{yes}.
\end{quote}
That is, if we knew the answer to $r(x)$, it would tell us the answer to $x$.

And the reduction $r$ is a \emph{polynomial time} reduction
if there is an algorithm that, given any input $x$ outputs $r(x)$,
with worst-case run time bounded by a polynomial in the length of $x$.
The notation
\[
  A \le_P B
\]
means ``$A$ reduces in polynomial time to $B$''.
It means that (up to polynomial factors)
$A$ is ``no harder'' than $B$:
if $A \le_P B$, and $B$ has a polynomial-time algorithm, then so does $A$.
This is because composing the reduction with any polynomial-time algorithm for $B$
gives a polynomial-time algorithm for $A$.

\newpage 

\subsection{Example 1: \prob{Bipartite Matching} $\le_P$ \prob{Max Flow}}
To illustrate the concept, consider the following decision problem:

\begin{mylist}
\item[] \prob{Bipartite Matching} (decision version) \dotfill \em
% [review 2026-09-27] fix: the bound was called ell here but k in the reduction and its proofs
\item[\bf input:] a bipartite graph $G=(U, W ,E)$, and a value $k$
\item[\bf output:] `\textsf{yes}' if $G$ has matching of size at least $k$, and otherwise `\textsf{no}'
\end{mylist}

Note that a \emph{bipartite} graph $G=(U, W, E)$
is an undirected graph in which the vertices are divided into two parts (sets)  $U$ and $W$,
and each edge $(u,w)$ in $E$ is in $U\times W$ (that is, $u$ is in $U$ and $w$ is in $W$, so all edges
go between the parts---no edge has both endpoints in one part).

\emph{[Consider an example for intuition.]}

Next we define a polynomial-time reduction $r$ from this problem to \prob{Max Flow}.
To define the reduction, one specifies, for every possible input $I$ (for \prob{Bipartite Matching})
what the output $r(I)$ is.  The output must be some input for \prob{Max Flow}.
We have to explicitly define this input in terms of $I$ (that is, from the bipartite graph $G$ and $k$).
Here is the definition of our reduction $r$:

\paragraph{Step 1.  Define the reduction.} Do this by specifying its output for every possible input.

\begin{mylist}
\item[] Our reduction $r$ from \prob{Bipartite Matching} to \prob{Max Flow}: \em
\item[\bf input:] any \prob{Bipartite Matching} instance $I=(G=(U, W ,E), k)$
\item[\bf output:] the output $r(I)$ is the \prob{Max Flow} instance $(G'=(V', E'), s, t, k')$ where
  \begin{align*}
    k' & {} = k \\
    % [review 2026-09-27] fix: V was undefined for a bipartite graph G=(U,W,E)
    V' & {} = U\cup W\cup\{s, t\} \\
    E' & {} = E \cup \{(s, u) : u\in U\} \cup \{(w, t) : w\in W\}
  \end{align*}
  and all edges have capacity 1.
\end{mylist}

In words: The reduction takes $k' = k$, then obtains $G'$ from $G$ as follows.
The vertices of $G'$ are the vertices of $G$ and two new vertices $s$ and $t$.
For each edge $(u, w)$ from $U$ to  $W$ in $G$,
there is a directed edge $(u, w)$ in $G'$ with capacity 1.
For each vertex $u\in U$,
there is a directed edge $(s, u)$ in $G'$ with capacity 1. 
For each vertex $w\in W$,
there is a directed edge $(w, t)$ in $G'$ with capacity 1. 

\smallskip

\emph{[Consider an example for intuition.  See also \KT Section 7.5, which presents the same reduction.]}

\paragraph{Step 2.  Explain how the reduction can be computed in polynomial time.}
To know that the reduction $r$ is a \emph{polynomial-time} reduction,
we have to be sure that, given any input $I$, the output $r(I)$
can be computed in time polynomial in the size of $I$.

By inspection of the reduction $r$ defined above, given $G$ and $k$, 
$G'$ (and $k'$) can be computed in time linear in the size of $G$,
by simply enumerating the vertices and edges of $G$ appropriately.
So the reduction can be computed in polynomial (even linear) time.

\paragraph{Step 3.  Prove that the reduction is correct.}
To know that the reduction $r$ is \emph{correct},
we have to be sure that, for every possible input $I$, 
the answer to the output $r(I)$ is \textsf{yes}
if and only if the answer to $I$ is \textsf{yes}.
(In other words, to solve $I$, it suffices to solve $r(I)$.)

Next we'll give a detailed proof.
Study the form of this proof carefully.
The reasoning for the correctness of any reduction
always has the same overall structure.
Only the details vary, depending on the specific reduction.

Note: Don't be intimidated by the length of the proof.
Most of the steps in the proof are straightforward.
Make sure you understand the high-level structure of the reasoning.
A short-form proof follows the long-form proof.  You may wish to review it first.

\begin{lemma}
  The reduction $r$ is correct.  That is, for every instance $I$ of \prob{Bipartite Matching},
  the answer to $I$ is \textsf{yes} if and only the answer to $r(I)$ (as an instance of \prob{Max Flow})
  is  \textsf{yes}.
\end{lemma}
\begin{longFormProof}
  \step Consider any possible input $I=(G=(U, W, E), k)$.  \comment{block: $\forall$}
  \step Let $G'=(V', E')$ and $k'$ be the output $r(I)$ of the reduction given $I$.
  \step First, we'll prove that, if the answer to $I$ is \textsf{yes},
  then the answer to $r(I)$ is \textsf{yes}.
  \\
  Then, we'll prove the converse: 
  if the answer to $r(I)$ is \textsf{yes},
  then the answer to $I$ is \textsf{yes}.
  \smallskip
  \begin{block}[flow: IF]
    \step Suppose that the answer to $I$ is \textsf{yes}.  \comment{block: if \ldots~then \ldots}
    \step That is, $G$ has some matching, say $M$, of size at least $k$.
    \step We'll show that $G'$ has some $s$-$t$ flow of size at least $k'=k$.
    \step That is, we'll show that the answer to $r(I)$ is \textsf{yes}. 
    \step Define flow $f$ as follows:
    \\~~~for every edge $(u,w)$ in $M$,
    let $f(s, u) = 1$,
    let $f(u, w) = 1$,
    and 
    let $f(w, t) = 1$.
    \\ (That is, send 1 unit of flow along the path $(s, u, w, t)$.)
    \\ For every other edge $(u, w)$ in $G'$, let $f(u, w) = 0$.
    \step By inspection of $f$, the value of $f$ is the size of $M$, which is at least $k'=k$.
    \step And $f$ satisfies the conservation and capacity constraints, so $f$ is a valid flow.
    \step So $G'$ has an $s$-$t$ flow of size at least $k'$.
    \step So, the answer to $r(I)$ is \textsf{yes}.
  \end{block}
  \step[flow: IF2] By Block~\ref{flow: IF}, if the answer to $I$ is \textsf{yes}, then the answer to $r(I)$ is \textsf{yes}.
  \medskip
  \step Next we show the converse.
  \begin{block}[flow: ONLYIF]
    \step Suppose that the answer to $r(I)$ is \textsf{yes}. \comment{block: if \ldots~then \ldots}
    \step That is, $G'$ has some $s$-$t$ flow, of value at least $k'=k$.
    \step Since $G'$ has integer capacities, $G'$ has such a flow for which each edge flow is an integer.
    \step Let $f$ be such a flow.
     Since each edge capacity in $G'$ is 1, each edge flow $f(u,w)$ is either 0 or 1.
    \step We'll show that $G$ has some matching of size at least $k$.
    \step Define $M = \{(u,w)\in E : f(u,w) = 1\}$
    to contain each edge in $G$ that has flow in $G'$.
    \step For each vertex $u\in U$, there is only one edge into $u$ in $G'$.
    \step So, because of the conservation constraints, each vertex $u\in U$ is in at most one edge in $M$. 
    \step Likewise, each vertex $w\in W$ is in at most one edge in $M$.
    \step So $M$ is a matching in $G$.
    \step Also by the conservation constraints,
    % [review 2026-09-27] fix: u' was a typo for u
    for each edge $(s, u)$ in $G'$ with $f(s, u) = 1$
    there is some edge $(u,w)$ in $M$,
    so the size of $M$ is at least the value of $f$, that is, at least $k$.
    \step So $M$ is a matching of size at least $k$ in $G$.
    \step So the answer to $I$ is \textsf{yes}. 
  \end{block}
  
  \medskip
  
  \step By Block~\ref{flow: ONLYIF}, if the answer to $r(I)$ is \textsf{yes}, then the answer to $I$ is \textsf{yes}.
  \step By this and Step~\ref{flow: IF2},
  the answer to $I$ is \textsf{yes} if and only if the answer to $r(I)$ is \textsf{yes}.
\end{longFormProof}

As mentioned above, the high-level structure of the proof of correctness is the same
for all reductions.  Only the details (within the two blocks) depend on the specific reduction.
For readers who understand this, it is enough to give just those details.
We could do that as follows:

\newpage

\begin{proof}[Proof of correctness, short form.]
  Consider any possible input $I=(G=(U,W,E), k)$ to the reduction.
  Let $G'=(V', E')$ and $k'=k$ be the output of the reduction for that input.

  Suppose that $G$ has some matching $M$ of size at least $k$.
  Define a flow $f$ in $G'$ by taking $f(s, u) = f(u, w) = f(w, t) = 1$
  for every matched edge $(u, w)$, and giving zero flow to the remaining edges.
  Then $f$ is an $s$-$t$ flow whose value is equal to the size of $M$, which is at least $k$.
  So, if $G$ has a matching of size at least $k$, then $G'$ has a flow of value at least $k'$.

  Conversely, suppose that $G'$ has a flow of value at least $k'=k$.
  Let $f$ be such a flow that is also integer-valued
  (there must be one, since the capacities in $G'$ are integers).
  So, the flow on each edge is either 1 or 0.
  Let $M$ contain the edges in $E$ whose corresponding edges in $G'$ have flow 1.
  For each vertex $u\in U$, there is only one edge into $u$ in $G'$,
  so (by the conservation constraints), each vertex $u\in U$ is in at most one edge in $M$.
  Likewise, each vertex $w\in W$ is in at most one edge in $M$. 
  So $M$ is a matching in $G$.
  Since the value of $f$ is at least $k$,
  there are at least $k$ edges out of $s$ that have flow 1.
  So, by inspection of $G'$ (and the conservation constraints),
  there are at least $k$ edges in $M$.
  So $M$ is a matching of size at least $k$ in $G$. 

  By the above reasoning, $G$ has a matching of size at least $k$
  if and only if $G'$ has an $s$-$t$ flow of value at least $k'$.
  That is, the answer to $I$ is \textsf{yes}
  if and only if the answer to $r(I)$ is \textsf{yes}.
  So the reduction is correct.
\end{proof}

\subsection{Example 2: \prob{Disjoint Paths} $\le_P$ \prob{Max Flow}} 
We give one more reduction, from the following problem to \prob{Max Flow}.

\begin{mylist}
\item[] \prob{Disjoint Paths} (decision version) \dotfill \em
\item[\bf input:] a digraph $G=(V ,E)$, source $s$, sink $t$, and a value $k$
\item[\bf output:] `\textsf{yes}' if $G$ there are $k$ $s$-$t$ paths in $G$ such that
  no two share an edge; otherwise `\textsf{no}'.
\end{mylist}

\emph{[Consider an example for intuition.]}

Next we define a polynomial-time reduction $r$ from this problem to \prob{Max Flow}.
To define the reduction, we specify, for every possible input $I$ (for \prob{Disjoint Paths})
what the output $r(I)$ is.  The output must be some input for \prob{Max Flow}, that is, a flow network and a value $k$.
We have to say specifically how to obtain this flow network from $I$ (that is, from $G$, $s$, $t$, and $k$).
Here is the definition of our reduction $r$:

\paragraph{Step 1.  Define the reduction.} (Specify the output for every possible input.)

\begin{mylist}
\item[] Here is our reduction $r$ from \prob{Disjoint Paths} to \prob{Max Flow}:
\item[\bf input:] any \prob{Disjoint Paths} instance $I=(G=(V ,E), s, t, k)$
\item[\bf output:] the output $r(I)$ is the \prob{Max Flow} instance $(G', s, t, k)$ where
  $G'$ is obtained from $G$ by assigning capacity 1 to each edge.
\end{mylist}

\emph{[Consider an example for intuition.  See also \KT Section 7.6, which presents the same reduction.]}

\paragraph{Step 2.  Explain how the reduction can be computed in polynomial time.}
By inspection of the reduction $r$ (as defined above), given $I=(G, s, t, k)$, 
the flow network $G'$ can be computed in time linear in the size of $G$,
by simply adding a capacity to each edge.
So the reduction can be computed in polynomial (even linear) time.

\paragraph{Step 3.  Prove that the reduction is correct.}
Following the previous reduction from \prob{Bipartite Matching},
we'll assume the reader understands the general high-level structure
common to all proofs of correctness for reductions.
We'll give a short-form proof. Before that, we need a utility lemma.
\begin{lemma}\label{flow: utility2}
  Given any $s$-$t$ flow $f$ with positive value,
  $t$ is reachable from $s$ using only edges that have positive flow.
\end{lemma}
\begin{shortFormProof}
  Let $G''$ be the flow network obtained by reducing the capacity of each edge to equal the flow on the edge,
  then deleting all edges that have flow (and modified capacity) zero.
  The original flow remains valid in $G''$.
  
  Suppose for contradiction that there is no $s$-$t$ path in $G''$.
  Let $S$ denote the set of vertices reachable from $s$ in $G''$,
  so that $(S, T=V\setminus S)$ is an $s$-$t$ cut.
  Note that no edge $(u,w)$ in $G''$ leaves $S$ (otherwise $u$ would be reachable from $s$,
  but $w$ would not be, even though there is an edge from $u$ to $w$).
  So, $(S, T)$ is an $s$-$t$ cut of capacity zero in $G''$.
  So, every $s$-$t$ flow in $G''$ has value at most zero.
  This is a contradiction, as the original flow is valid in $G''$ and has positive value.
\end{shortFormProof}

Here is the proof that the reduction is correct.
\begin{lemma}
  The reduction $r$ is correct.  That is, for every input $I$,
  the answer to $I$ is \textsf{yes} if and only if the answer to $r(I)$ is \textsf{yes}.
\end{lemma}
\begin{shortFormProof}
  Consider any possible input $I=(G=(V, E), s, t, k)$ to the reduction.
  Let $G'$ (and $s, t, k$) be the output of the reduction for that input.

  First suppose that $G$ does have $k$ edge-disjoint $s$-$t$ paths,
  say, $P_1, P_2, \ldots, P_k$.
  Assume without loss of generality that each path is simple.
  (If not, just remove all cycles from each path to get $k$ simple edge-disjoint paths.)
  Define a flow $f$ in $G'$ by setting $f(u, w)$ to be the number
  of times $(u,w)$ occurs on these paths.
  Since the paths are edge-disjoint and simple,
  each flow value $f(u,w)$ will be 0 or 1.
  Since each edge capacity  in $G'$ is 1, $f$ obeys the capacity constraints.
  Since each of the paths leaves $s$,
  the value of $f$ is $k$.
  For each vertex $v$ (other than $s$ or $t$),
  for each edge $(u, v)$ on some path $P_i$ entering $v$,
  the next edge $(v, w)$ on $P_i$ leaves $v$,
  so the total flow into $v$ equals the total flow leaving $v$.
  So $f$ obeys the conservation constraints.
  So $f$ is a valid flow in $G'$, of value $k$.

  Conversely, suppose that $G'$ does have a flow of value at least $k$.
  All edge capacities are integers, 
  so there is also an \emph{integer} flow of value at least $k$.
  Let $f$ be such a flow.
  Since each edge capacity is 1, each flow value $f(u,w)$ is 0 or 1.
  Convert $f$ into a sequence of paths $P_1, P_2, \ldots, P_\ell$ as follows:
  \begin{myframed}
    \begin{steps}
      \step for $i= 1, 2, \ldots$, while the value of flow $f$ is positive:
      \begin{block}
        \step let $P_i$ be an $s$-$t$ path that uses only edges $(u,w)$ such that $f(u,w)=1$
        \step for each edge $(u,w)\in P_i$, reduce the flow $f(u,w)$ on $(u,w)$ by 1 (to 0)
      \end{block}
    \end{steps}
  \end{myframed}
  In each iteration, the $s$-$t$ path $P_i$ must exist by Lemma~\ref{flow: utility2}.
  The iteration reduces the flow on each edge in $P_i$ by 1 (to 0).
  This preserves the conservation constraints, the capacity constraints,
  and the property that the flow on each edge is either 0 or 1.
  But it reduces the value of the flow by 1.
  It follows that the number of iterations will equal the original value of $f$.
  Since this is at least $k$, the process produces at least $k$ paths.
  And no edge occurs twice in these paths,
  because an edge $(u,w)$ is used in a path $P_i$ only when $f(u,w)$ is 1,
  and after that $f(u,w)$ is reduced to 0.
  So, if $G'$ has a flow of value at least $k$,
  then $G$ has $k$ edge-disjoint $s$-$t$ paths.
  (If the process produces more than $k$ paths, just ignore some.)
\end{shortFormProof}

\subsection{Exercise: \prob{Bipartite Matching} $\le_P$ \prob{Disjoint Paths}}

As an exercise, find a poly-time reduction from \prob{Bipartite Matching} to \prob{Disjoint Paths}.

\subsection{The two key properties of polynomial-time reductions} 
Recall that $A \le_P B$ denotes that
there is a polynomial-time reduction from decision problem $A$
to decision problem $B$.
Intuitively, this notation suggests that (up to polynomial time),
$A$ is ``no harder'' than $B$.
Formally, the following holds:

\begin{lemma}\label{flow: reduction}
  Let $A$ and $B$ be decision problems
  such that $A$ reduces in polynomial time to $B$.
  If $B$ has a polynomial-time algorithm,
  then $A$ has a polynomial-time algorithm.
  That is, if $A \le_P B$ and $B\in P$, then $A\in P$.
\end{lemma}
\begin{proof}
  Assume $r$ is a poly-time reduction from $A$ to $B$,
  and that $B$ has a poly-time algorithm.
  Fix $c$ and $c'$ such that $r$ is computable in time $O(n^c)$
  and $B$ has an algorithm running in time $O(n^{c'})$.
  Consider the following algorithm for $A$:
  \begin{myframed}
    \begin{steps}
    \item[] algorithm for $A$ (on input $x$):
      \step compute $y= r(x)$, in time $O(|x|^c)$
      \step run the algorithm
      for $B$ on input $y$, taking time $O(|y|^{c'})$
      \step if the algorithm for $B$ outputs \textsf{yes}, then output \textsf{yes}, otherwise \textsf{no}
    \end{steps}
  \end{myframed}
  \smallskip
  
  This algorithm for $A$ is correct by the following reasoning:

  The algorithm for $A$ outputs \textsf{yes} on input $x$

  $\iff$ the algorithm for $B$ outputs \textsf{yes} on input $y$
  \comment{(by inspection of the algorithm for $A$)}

  $\iff$ the answer for $r(x)$ is \textsf{yes}
  \comment{(as the algorithm for $B$ is correct)}

  $\iff$ the answer for $x$ is \textsf{yes}
  \comment{(because $r$ is a correct reduction from $A$ to $B$)}

  So the algorithm is correct.
  
  The algorithm takes polynomial time, $O(|x|^{c\,c'})$, by the following reasoning.
  Step 1, computing $y=r(x)$, takes time $O(|x|^c)$
  because $r(x)$ is computable in time $O(|x|^c)$.
  Also, this means that $y$ has size $O(|x|^c)$.
  So Step 2, running the algorithm for $B$ on $y$,
  takes time $O(|y|^{c'}) = O( (|x|^c)^{c'}) = O(|x|^{c\,c'})$.
\end{proof}

One more formal observation:
poly-time reductions compose,
so the relation $\le_P$ is transitive:

\begin{lemma}\label{flow: transitive}
  Let $A$, $B$, and $C$ be decision problems.
  If $A$ reduces in polynomial time to $B$,
  and $B$ reduces in polynomial time to $C$,
  then $A$ reduces in polynomial time to $C$.

  That is, if $A\le_P B$ and $B\le_P C$, then $A\le_P C$.
\end{lemma}
\begin{proof}
  Let $r_1$ and $r_2$ be the poly-time reductions from $A$ to $B$, and from $B$ to $C$.
  Consider the function $r_3$ defined by composing $r_1$ and $r_2$,
  that is, $r_3(x) = r_2(r_1(x))$.
  The function $r_3$ is a correct reduction from $A$ to $C$ by the following reasoning:

  The answer for $x$ is \textsf{yes}

  $\iff$ the answer for $r_1(x)$ is \textsf{yes}
  \comment{(because $r_1$ reduces $A$ to $B$)}

  $\iff$ the answer for $r_2(r_1(x))$ is \textsf{yes}
  \comment{(because $r_2$ reduces $B$ to $C$)}

  $\iff$ the answer for $r_3(x)$ is \textsf{yes}
  \comment{(because $r_3(x) = r_2(r_1(x))$)}

  \smallskip
  So the reduction $r_3$ is correct.
  
  For the running time, fix $c_1$ and $c_2$ such that
  $r_1$ is computable in time $O(n^{c_1})$
  and $r_2$ is computable in time $O(n^{c_2})$.
  The following algorithm computes $r_3$ in polynomial time, $O(n^{c_1 c_2})$:

  \begin{myframed}
    \begin{steps}
    \item[] algorithm to compute $r_3(x)$:
      \step compute $y = r_1(x)$, in time $O(|x|^{c_1})$
      \step compute and return $r_2(y)$, taking time $O(|y|^{c_2})$
    \end{steps}
  \end{myframed}
  \smallskip

  Clearly the algorithm computes $r_3(x)$.
  To bound the run time,
  reason as in the previous proof:
  Step 1 takes time $O(|x|^{c_1})$,
  so $y$ has size $|y|=O(|x|^{c_1})$,
  so Step 2 takes time $O(|y|^{c_2}) = O(|x|^{c_1 c_2})$.
\end{proof}

\subsection{External sources on polynomial-time reductions}


\begin{itemize}
\item Reducing matching to flow:
  \KT Section 7.5;
  \DPV Section 7.3;
  \CLRS Section 26.3

\item Reducing disjoint paths to flow:
  \KT Section 7.6

\item For \prob{Seam Carving} (which reduces to shortest $s$-$t$ path in a DAG), see
  
  \URLplain{https://en.wikipedia.org/wiki/Seam_carving}

  \URLplain{https://www.youtube.com/watch?v=6NcIJXTlugc}

\item For \prob{Image Segmentation} via \prob{Min Cut}, see \KT Section 7.10.  See also:

  \url{https://discovery.ucl.ac.uk/13383/1/13383.pdf}, Figures 1, 2, and 11

\item Other reductions to \prob{Max Flow} or \prob{Min Cut}:

  \KT Sections 7.5--7.12 (and \KWSlides for Ch.~7); see also \JE Chapter 11 

\item Reductions in general:

  \KT Section 8.1 (and \KWSlides for Ch.~8);  \CLRS Section 34.3; and \JE Chapter 12, and
  \URLplain{https://ocw.mit.edu/courses/electrical-engineering-and-computer-science/6-006-introduction-to-algorithms-fall-2011/lecture-videos/lecture-23-computational-complexity/}

\end{itemize}

%%% Local Variables:
%%% mode: latex
%%% TeX-master: "flow_reductions_LP"
%%% End:
